\documentclass[11pt,a4paper]{article}
\usepackage[margin=1in]{geometry}
\usepackage[T1]{fontenc}
\usepackage[utf8]{inputenc}
\usepackage{lmodern}
\usepackage{hyperref}
\usepackage{amsmath,amssymb}

\title{Informed Machine Learning Notebook}
\author{}
\date{\today}

\begin{document}
\maketitle

\section{Overview about these five papers}
These five papers show me the overview about informed machine learning:
\begin{itemize}
\item taxonomy-level Informed ML,
\item logic-constrained neural learning,
\item granular computing for confidence-aware outputs.
\end{itemize}

\section{A Unified Optimization Template}
A compact mathematical template that covers most methods in this corpus is:
\begin{equation}
\min_{\theta} \; \mathcal{L}_{\text{task}}(\theta)
+ \lambda_{K}\,\mathcal{L}_{\text{knowledge}}(\theta)
+ \lambda_{G}\,\mathcal{L}_{\text{granularity}}(\theta).
\label{eq:unified}
\end{equation}

Interpretation:
\begin{itemize}
\item $\mathcal{L}_{\text{task}}$: standard supervised objective (e.g., cross-entropy, MSE).
\item $\mathcal{L}_{\text{knowledge}}$: prior-knowledge term (logic rules, equations, graph constraints).
\item $\mathcal{L}_{\text{granularity}}$: uncertainty/credibility term that promotes informative non-point outputs.
\end{itemize}
Equation \eqref{eq:unified} is the mathematical counterpart of the taxonomy path ``knowledge source $\rightarrow$ representation $\rightarrow$ integration location'' in the TKDE survey.

\section{Logic-Constrained Learning: Three Mathematical Mechanisms}

\subsection{Hu et al. (logic-net): Posterior Regularization + Distillation}
Let $p_{\theta}(y\mid x)$ be the student network. A rule-regularized teacher $q(y\mid x)$ is constructed by projection:
\begin{equation}
q^{\ast} = \arg\min_{q \in \mathcal{Q}} \mathrm{KL}\big(q(Y\mid X)\,\|\,p_{\theta}(Y\mid X)\big) 
+ C \sum_{\ell} \xi_{\ell},
\label{eq:pr_projection}
\end{equation}
subject to rule satisfaction constraints (soft form) using slack variables $\xi_{\ell}$.

A common closed-form structure is:
\begin{equation}
q^{\ast}(y\mid x) \propto p_{\theta}(y\mid x)
\exp\!\left\{-\sum_{\ell} \gamma_{\ell}\,\big(1-r_{\ell}(x,y)\big)\right\},
\label{eq:teacher}
\end{equation}
where $r_{\ell}(x,y) \in [0,1]$ is soft truth of rule $\ell$.

Then student update mixes data labels and teacher supervision:
\begin{equation}
\theta \leftarrow \arg\min_{\theta} \frac{1}{N}\sum_{n=1}^{N}
\Big[(1-\pi_t)\,\ell\big(y_n,p_{\theta}(\cdot\mid x_n)\big)
+ \pi_t\,\ell\big(q^{\ast}(\cdot\mid x_n),p_{\theta}(\cdot\mid x_n)\big)\Big].
\label{eq:distill}
\end{equation}
This is the ``teacher-student + rule projection'' logic in mathematical form.

\subsection{Xu et al. (semantic loss): Constraint-Exact Differentiable Loss}
For Boolean outputs $X_1,\dots,X_n$ with network probabilities $p=(p_1,\dots,p_n)$ and a logical constraint $\alpha$, semantic loss is:
\begin{equation}
\mathcal{L}_{s}(\alpha,p)
= -\log \sum_{x \models \alpha}
\prod_{i: x_i=1} p_i \prod_{i: x_i=0}(1-p_i).
\label{eq:semantic_loss}
\end{equation}

This directly measures probability mass assigned to satisfying worlds of $\alpha$. Training usually uses:
\begin{equation}
\mathcal{L} = \mathcal{L}_{\text{task}} + \beta\,\mathcal{L}_{s}(\alpha,p).
\label{eq:semantic_total}
\end{equation}
Key property: logical meaning is preserved exactly at the constraint level, while optimization remains gradient-based.

\subsection{Fischer et al. (DL2): Declarative Constraints to Loss and Queries}
DL2 maps constraints $\varphi$ into nonnegative losses $\mathcal{L}(\varphi)$ with:
\begin{equation}
\mathcal{L}(\varphi)=0 \iff \varphi \text{ is satisfied},
\qquad
\mathcal{L}(\varphi) \ge 0,
\label{eq:dl2_property}
\end{equation}
and differentiability almost everywhere.

Representative atom relaxations (hinge-style) are:
\begin{equation}
\mathcal{L}(a \le b)=\max(0,a-b), \qquad
\mathcal{L}(a=b)=|a-b|.
\label{eq:atom_relax}
\end{equation}
Complex formulas are recursively composed.

Training with constraints can be written as:
\begin{equation}
\min_{\theta} \; \mathcal{L}_{\text{task}}(\theta)
+ \lambda\,\mathbb{E}_{x}\big[\mathcal{L}(\varphi(x,\theta))\big].
\label{eq:dl2_train}
\end{equation}

Querying is posed as input optimization:
\begin{equation}
x^{\ast}=\arg\min_{x}\;\mathcal{L}(\psi(x,\theta)),
\label{eq:dl2_query}
\end{equation}
where $\psi$ encodes the query constraints.

\section{From Point Predictions to Granular Predictions}

\subsection{Takagi--Sugeno Base Model (numeric)}
For rule $i$:
\begin{equation}
R_i:\; \text{if }x_1\text{ is }A_{i1},\dots,x_m\text{ is }A_{im},\text{ then } y=b_i,
\label{eq:ts_rule}
\end{equation}
with firing strength
\begin{equation}
w_i(x)=\prod_{j=1}^{m}\mu_{A_{ij}}(x_j).
\label{eq:firing}
\end{equation}
Output:
\begin{equation}
\hat y(x)=\frac{\sum_{i=1}^{c} w_i(x)b_i}{\sum_{i=1}^{c}w_i(x)},
\quad
\mathcal{L}_{\text{task}}=\sum_{k}(\hat y(x_k)-y_k)^2.
\label{eq:ts_output}
\end{equation}

\subsection{Granular Elevation}
Instead of scalar consequents $b_i$, use granules $B_i$ (interval/fuzzy/probabilistic):
\begin{equation}
B_i = [\underline b_i,\overline b_i] \quad \text{(interval case)}.
\label{eq:interval_consequent}
\end{equation}
Then model output becomes a granular object, e.g., interval:
\begin{equation}
\hat Y(x)=\left[
\frac{\sum_i w_i(x)\underline b_i}{\sum_i w_i(x)},
\frac{\sum_i w_i(x)\overline b_i}{\sum_i w_i(x)}
\right].
\label{eq:interval_output}
\end{equation}
This converts ``single value'' prediction into ``value + credibility span'' prediction.

\subsection{Principle of Justifiable Granularity}
A granule should balance:
\begin{itemize}
\item \textbf{coverage} $\mathrm{cov}(G)$: how much data are represented,
\item \textbf{specificity} $\mathrm{sp}(G)$: how informative/narrow the granule is.
\end{itemize}
A standard construction objective is:
\begin{equation}
G^{\ast}=\arg\max_{G}\;V(G),
\qquad
V(G)=\mathrm{cov}(G)\cdot \mathrm{sp}(G).
\label{eq:jg}
\end{equation}
This is the key mathematical principle used to build interval/fuzzy granules in the 2025 TFS paper.

\subsection{Probabilistic Granules via Gaussian Process}
For GP regression:
\begin{equation}
y(x)\mid \mathcal{D} \sim \mathcal{N}(\mu(x),\sigma^2(x)).
\label{eq:gp_pred}
\end{equation}
A $(1-\alpha)$ confidence interval is:
\begin{equation}
\hat Y_{\alpha}(x)=\big[\mu(x)-z_{\alpha/2}\sigma(x),\;\mu(x)+z_{\alpha/2}\sigma(x)\big].
\label{eq:gp_ci}
\end{equation}
The paper further maps probabilistic granules to interval/fuzzy granules through justifiable granularity criteria.

\section{A Deeper Synthesis: What Each Family Optimizes}
\begin{itemize}
\item \textbf{logic-net}: optimize predictive fit while distilling rule-regularized posterior structure.
\item \textbf{semantic loss}: optimize probability mass over logically valid worlds.
\item \textbf{DL2}: optimize satisfiability of declarative constraints in both training and querying.
\item \textbf{granular TS/GC}: optimize not only fit, but also the informativeness of uncertainty representation.
\end{itemize}

So, the progression is:
\begin{equation}
\text{validity of output} \;\Rightarrow\; \text{validity + confidence geometry of output}.
\label{eq:progression}
\end{equation}

\section{A Research Blueprint}
A mathematically strong formulation for my own work can be:
\begin{equation}
\min_{\theta,\,\eta}\;
\underbrace{\mathcal{L}_{\text{task}}(\theta)}_{\text{fit}}
+ \lambda_1\underbrace{\mathcal{L}_{\text{logic}}(\theta)}_{\text{rule satisfaction}}
+ \lambda_2\underbrace{\mathcal{L}_{\text{gran}}(\theta,\eta)}_{\text{credible granules}}
+ \lambda_3\underbrace{\Omega(\theta,\eta)}_{\text{regularization}},
\label{eq:blueprint}
\end{equation}
where $\eta$ parameterizes interval/fuzzy/probabilistic output structure.

Recommended evaluation should report at least:
\begin{itemize}
\item prediction quality (Acc/F1/RMSE),
\item rule satisfaction rate,
\item granule quality (coverage, specificity, and calibration metrics).
\end{itemize}


\end{document}
